\chapter{Reproducción}
\label{app:reproduccion}

Este apéndice describe el procedimiento utilizado para regenerar datos, entrenar las CNN, construir los informes y compilar la memoria. Los comandos se ejecutan desde la raíz del repositorio.

\section{Preparación del entorno}

\begin{lstlisting}[language=bash]
uv sync --extra dev --extra cnn --extra real-data
uv run --no-sync python --version
\end{lstlisting}

Para las campañas CNN se recomienda fijar el backend de Matplotlib:

\begin{lstlisting}[language=bash]
export MPLBACKEND=Agg
\end{lstlisting}

\section{Generación de datos}

El conjunto sintético QRS/ST se regenera con:

\begin{lstlisting}[language=bash]
scripts/run/build_simulator_qrs_dataset.sh
\end{lstlisting}

El conjunto Brugada HUCA procesado se construye desde la descarga local de PhysioNet:

\begin{lstlisting}[language=bash]
scripts/run/build_brugada_huca_dataset.sh
\end{lstlisting}

Los resultados esperados son:

\begin{itemize}
  \item \texttt{data/simulator\_qrs}: 100 pacientes, 300 imágenes y pliegues LOOCV;
  \item \texttt{data/brugada\_huca}: 317 pacientes incluidos, 951 imágenes y 46 exclusiones documentadas.
\end{itemize}

\section{Entrenamiento CNN}

La campaña sintética se lanza con:

\begin{lstlisting}[language=bash]
RESUME=0 scripts/run/run_cnn_simulator_qrs_loocv.sh
\end{lstlisting}

La campaña real HUCA se lanza con:

\begin{lstlisting}[language=bash]
RESUME=0 scripts/run/run_cnn_loocv.sh
\end{lstlisting}

Ambos scripts usan por defecto \(K=2,4,8,16,32\), semillas 42, 123 y 2026, ResNet18, 20 épocas, lote 32 y selección de umbral por exactitud equilibrada de validación.

\section{Adaptación de dominio}

La rama de adaptación se ejecuta método a método:

\begin{lstlisting}[language=bash]
METHOD=ssl   RESUME=0 scripts/run/run_cnn_domain_adaptation_loocv.sh
METHOD=coral RESUME=0 scripts/run/run_cnn_domain_adaptation_loocv.sh
METHOD=mmd   RESUME=0 scripts/run/run_cnn_domain_adaptation_loocv.sh
METHOD=dann  RESUME=0 scripts/run/run_cnn_domain_adaptation_loocv.sh
\end{lstlisting}

El método \texttt{ssl} se reproduce como preentrenamiento SimCLR de tres épocas y ajuste supervisado sin pérdida adicional de adaptación. CORAL, MMD y DANN usan el peso de adaptación 0,1. La salida queda en \texttt{outputs/cnn\_domain\_adaptation/} y \texttt{reports/loocv/cnn\_domain\_adaptation/}.

\section{Comparaciones e informes}

La comparación sintético-real se reconstruye con:

\begin{lstlisting}[language=bash]
scripts/run/build_loocv_comparison.sh
\end{lstlisting}

La comparación de adaptación de dominio se genera con:

\begin{lstlisting}[language=bash]
uv run --no-sync python scripts/eval/compare_cnn_domain_adaptation_reports.py
\end{lstlisting}

La auditoría final se lanza con:

\begin{lstlisting}[language=bash]
scripts/run/audit_loocv_results.sh
\end{lstlisting}

El cierre correcto debe incluir:

\begin{itemize}
  \item \texttt{all\_cnn\_experiments\_complete: OK};
  \item \texttt{cnn\_domain\_adaptation\_complete: OK};
  \item \texttt{final\_cnn\_package\_ready: OK};
  \item \texttt{vlm\_protocol\_ready: OK};
  \item \texttt{vlm\_numeric\_results: por\_ejecutar}.
\end{itemize}

\section{Ejecución VLM/ICL}

La infraestructura VLM/ICL queda preparada como segundo brazo de la comparación. La evaluación final debe lanzarse desde la raíz del repositorio con una única celda de comandos. El bloque valida primero el plan sin inferencia, ejecuta después el simulador QRS/ST, repite el proceso sobre Brugada-HUCA real usando el simulador como banco de demostraciones etiquetadas y, finalmente, reconstruye las comparaciones CNN frente a VLM.

\begin{lstlisting}[language=bash]
# 0) Ejecutar desde la raiz del repositorio.
set -eu
export MPLBACKEND=Agg
export VLM_API_BASE="${VLM_API_BASE:-http://servidor:8000/v1}"
export VLM_MODELS="${VLM_MODELS:-google/gemma-4-E4B-it,google/medgemma-4b-it}"
export VLM_RUNTIME="${VLM_RUNTIME:-remote_api}"
export K_VALUES="${K_VALUES:-0,2,4,8,16,32}"
export CONTROL_K_VALUES="${CONTROL_K_VALUES:-8,16,32}"
export SEEDS="${SEEDS:-42,123,2026}"
export CONDITIONS="${CONDITIONS:-zero_shot,normal,balanced,permuted,no_support_images}"

# 1) Validar plan VLM en el simulador QRS/ST, sin inferencia.
DATASET_ROOT="$PWD/data/simulator_qrs" \
CONTEXT_DATASET_ROOT="$PWD/data/simulator_qrs" \
REPORT_DIR="$PWD/reports/loocv/vlm_simulator_qrs" \
OUTPUT="$PWD/reports/loocv/vlm_simulator_qrs/vlm_setup_validation.json" \
scripts/run/validate_vlm_loocv.sh

# 2) Ejecutar inferencia VLM en el simulador QRS/ST.
DATASET_ROOT="$PWD/data/simulator_qrs" \
CONTEXT_DATASET_ROOT="$PWD/data/simulator_qrs" \
OUTPUT_ROOT="$PWD/outputs/vlm_simulator_qrs_loocv" \
REPORT_DIR="$PWD/reports/loocv/vlm_simulator_qrs" \
scripts/run/run_vlm_loocv.sh

# 3) Validar plan VLM en Brugada-HUCA real.
DATASET_ROOT="$PWD/data/brugada_huca" \
CONTEXT_DATASET_ROOT="$PWD/data/simulator_qrs" \
REPORT_DIR="$PWD/reports/loocv/vlm" \
OUTPUT="$PWD/reports/loocv/vlm/vlm_setup_validation.json" \
scripts/run/validate_vlm_loocv.sh

# 4) Ejecutar inferencia VLM en Brugada-HUCA real.
DATASET_ROOT="$PWD/data/brugada_huca" \
CONTEXT_DATASET_ROOT="$PWD/data/simulator_qrs" \
OUTPUT_ROOT="$PWD/outputs/vlm_loocv" \
REPORT_DIR="$PWD/reports/loocv/vlm" \
scripts/run/run_vlm_loocv.sh

# 5) Comparar CNN frente a VLM en simulador y en HUCA real.
CNN_SUMMARY="$PWD/reports/loocv/cnn_simulator_qrs/cnn_summary_by_seed.csv" \
VLM_SUMMARY="$PWD/reports/loocv/vlm_simulator_qrs/vlm_summary_by_seed.csv" \
VLM_CONDITION=normal \
OUTPUT_DIR="$PWD/reports/loocv/comparison_vlm_simulator_qrs" \
scripts/run/build_loocv_comparison.sh

CNN_SUMMARY="$PWD/reports/loocv/cnn/cnn_summary_by_seed.csv" \
VLM_SUMMARY="$PWD/reports/loocv/vlm/vlm_summary_by_seed.csv" \
VLM_CONDITION=normal \
OUTPUT_DIR="$PWD/reports/loocv/comparison" \
scripts/run/build_loocv_comparison.sh

# 6) Auditar que los artefactos finales existen.
scripts/run/audit_loocv_results.sh
\end{lstlisting}

Por defecto se ejecutan \(K=\{0,2,4,8,16,32\}\), semillas 42, 123 y 2026, y las condiciones \texttt{zero\_shot}, \texttt{normal}, \texttt{balanced}, \texttt{permuted} y \texttt{no\_support\_images}. Los controles usan \texttt{CONTROL\_K\_VALUES=8,16,32}. Si el servidor vLLM solo sirve un modelo cada vez, se cambia \texttt{VLM\_MODELS} a un único identificador y se repite la celda por modelo. Para probar la ruta de artefactos sin GPU ni API puede usarse:

\begin{lstlisting}[language=bash]
VLM_RUNTIME=local_gpu \
DRY_RUN_PREDICTIONS=negative \
LIMIT_FOLDS=2 \
K_VALUES=0,2 \
CONTROL_K_VALUES=2 \
MODELS=fake/gemma4,fake/medgemma \
scripts/run/run_vlm_loocv.sh
\end{lstlisting}

Antes de incorporar resultados VLM al documento deben existir, como mínimo:

\begin{itemize}
  \item \texttt{reports/loocv/vlm/vlm\_campaign\_manifest.json};
  \item \texttt{reports/loocv/vlm/vlm\_summary\_by\_seed.csv};
  \item \texttt{reports/loocv/vlm/vlm\_summary\_by\_k.csv};
  \item \texttt{reports/loocv/vlm/vlm\_summary\_by\_model\_condition\_k.csv};
  \item matrices de confusión en el subdirectorio \texttt{confusion\_by\_model\_condition\_k/} del informe VLM;
  \item gráficas \texttt{balanced\_accuracy\_by\_k.png} y \texttt{f1\_by\_k.png};
  \item predicciones por muestra bajo \texttt{outputs/vlm\_loocv/}, separadas por modelo, condición, \(K\) y semilla;
  \item auditoría con métricas VLM completas.
\end{itemize}

\section{Compilación de la memoria}

La memoria se compila desde el helper local de LaTeX:

\begin{lstlisting}[language=bash]
python scripts/compile_latex.py D:/Python_D/ecg/ecg-few/thesis/thesis/main.tex --compiler tectonic --json
\end{lstlisting}

Después de compilar, el PDF debe renderizarse a imágenes y revisarse visualmente para comprobar que las tablas, figuras, cabeceras, pies de página y referencias no presentan solapes ni elementos pendientes.
